\documentclass[10pt,a4paper]{amsart}
\usepackage{amsmath,amssymb,amsthm,graphicx}
% Retained minimal preamble for the Stacks Project chapter more-algebra.tex (More on Algebra).
% The chapter file pinned-original-source.tex carries only \input{preamble} on its first line, so no
% preamble text lives inside the chapter itself.  The Stacks Project uses one single shared preamble for
% every chapter; the snapshot stages it identically as sources/stacks-chow/preamble.tex and
% sources/stacks-etale-cohomology/preamble.tex.  The block below is copied mechanically, line for line,
% from sources/stacks-chow/preamble.tex lines 14-18 (the Xy-pic setup) and 154-217 (the native theorem
% declarations together with the mathematical macros).  Nothing was added, removed or redefined; the
% wrapper class and the amsmath/amssymb/amsthm/graphicx base packages are supplied by the assembler.
% Chapter: more-algebra.tex -- More on Algebra.
% For commutative diagrams we use Xy-pic
\usepackage[all]{xy}

% We use 2cell for 2-commutative diagrams.
\xyoption{2cell}

% Theorem environments.
%
\theoremstyle{plain}
\newtheorem{theorem}[subsection]{Theorem}
\newtheorem{proposition}[subsection]{Proposition}
\newtheorem{lemma}[subsection]{Lemma}

\theoremstyle{definition}
\newtheorem{definition}[subsection]{Definition}
\newtheorem{example}[subsection]{Example}
\newtheorem{exercise}[subsection]{Exercise}
\newtheorem{situation}[subsection]{Situation}

\theoremstyle{remark}
\newtheorem{remark}[subsection]{Remark}
\newtheorem{remarks}[subsection]{Remarks}

\numberwithin{equation}{subsection}

% Macros
%
\def\lim{\mathop{\mathrm{lim}}\nolimits}
\def\colim{\mathop{\mathrm{colim}}\nolimits}
\def\Spec{\mathop{\mathrm{Spec}}}
\def\Hom{\mathop{\mathrm{Hom}}\nolimits}
\def\Ext{\mathop{\mathrm{Ext}}\nolimits}
\def\SheafHom{\mathop{\mathcal{H}\!\mathit{om}}\nolimits}
\def\SheafExt{\mathop{\mathcal{E}\!\mathit{xt}}\nolimits}
\def\Sch{\mathit{Sch}}
\def\Mor{\mathop{\mathrm{Mor}}\nolimits}
\def\Ob{\mathop{\mathrm{Ob}}\nolimits}
\def\Sh{\mathop{\mathit{Sh}}\nolimits}
\def\NL{\mathop{N\!L}\nolimits}
\def\CH{\mathop{\mathrm{CH}}\nolimits}
\def\proetale{{pro\text{-}\acute{e}tale}}
\def\etale{{\acute{e}tale}}
\def\QCoh{\mathit{QCoh}}
\def\Ker{\mathop{\mathrm{Ker}}}
\def\Im{\mathop{\mathrm{Im}}}
\def\Coker{\mathop{\mathrm{Coker}}}
\def\Coim{\mathop{\mathrm{Coim}}}

% Boxtimes
%
\DeclareMathSymbol{\boxtimes}{\mathbin}{AMSa}{"02}

%
% Macros for moduli stacks/spaces
%
\def\QCohstack{\mathcal{QC}\!\mathit{oh}}
\def\Cohstack{\mathcal{C}\!\mathit{oh}}
\def\Spacesstack{\mathcal{S}\!\mathit{paces}}
\def\Quotfunctor{\mathrm{Quot}}
\def\Hilbfunctor{\mathrm{Hilb}}
\def\Curvesstack{\mathcal{C}\!\mathit{urves}}
\def\Polarizedstack{\mathcal{P}\!\mathit{olarized}}
\def\Complexesstack{\mathcal{C}\!\mathit{omplexes}}
% \Pic is the operator that assigns to X its picard group, usage \Pic(X)
% \Picardstack_{X/B} denotes the Picard stack of X over B
% \Picardfunctor_{X/B} denotes the Picard functor of X over B
\def\Pic{\mathop{\mathrm{Pic}}\nolimits}
\def\Picardstack{\mathcal{P}\!\mathit{ic}}
\def\Picardfunctor{\mathrm{Pic}}
\def\Deformationcategory{\mathcal{D}\!\mathit{ef}}

% --- appended 2026-10-12 (Chief): the project own definitions of the reference/slogan/history
% display environments, copied mechanically from the shared Stacks preamble lines 8-12
% (\usepackage{verbatim} plus the three \newenvironment lines, which the chapter preamble keeps
% verbatim from the project preamble). Nothing added, removed or redefined: these environments
% carry no mathematical content and the project itself typesets them as comments. Without these
% three lines a counted statement whose own text uses them fails with Environment slogan undefined.
\usepackage{verbatim}
\newenvironment{reference}{\comment}{\endcomment}
\newenvironment{slogan}{\comment}{\endcomment}
\newenvironment{history}{\comment}{\endcomment}
\usepackage{cleveref}
\crefname{originallemma}{original lemma}{original lemmas}
\Crefname{originallemma}{Original lemma}{Original lemmas}
\makeatletter
\newlabel{derived-section-cones}{{014D}{source}{}{sourcepointer-1}{}}
\newlabel{derived-section-cones@cref}{{[originallemma][0][]014D}{[1][0][]source}{}{}{}}
\makeatother
\crefname{originallemma}{original lemma}{original lemmas}
\Crefname{originallemma}{Original lemma}{Original lemmas}
\makeatletter
\newlabel{lemma-vanishing-coh-prod-totalization}{{30096}{source}{}{sourcepointer-2}{}}
\newlabel{lemma-vanishing-coh-prod-totalization@cref}{{[originallemma][0][]30096}{[1][0][]source}{}{}{}}
\makeatother
\title{A levelwise quasi-isomorphism of complexes of complexes induces a quasi-isomorphism of product totalizations}
\author{The Stacks Project Authors}
\date{stacks-project:more-algebra.tex:local-snapshot-20261012; GNU FDL 1.2 or later}
\begin{document}
\maketitle
\setcounter{subsection}{1}
\noindent\textit{Editorial note.} All delivered mathematical text is the author's own, copied from the pinned Stacks Project chapter file more-algebra.tex; no statement or proof step was rewritten, reordered or removed.\par\medskip
\noindent\textit{Editorial note.} The chapter file pinned-original-source.tex carries only \textbackslash{}input\{preamble\} on its first line, so no preamble text lives inside the chapter. The retained file sources/stacks-more-algebra/retained-preamble.tex is a mechanical, line-for-line copy of the project's single shared preamble as staged in the snapshot; nothing was added, removed or redefined and no package was added.\par\medskip
\noindent\textit{Editorial note.} The source URL is the source-level Stacks Project URL, as required when more than one card is produced from one source; the per-claim tag is carried by claim\_key.\par\medskip
\noindent\textit{Editorial note.} The author writes '(We omit the verification that the differentials agree.)' inside the counted proof span; that sentence is kept verbatim. The two references inside the counted proof span are recorded as non-load-bearing navigation pointers.\par\medskip
\noindent\textit{Allowed standard prerequisites (editorial).}\begin{itemize}
\item The cone construction on a map of complexes and its functoriality. invoked at proof lines 30152-30155. Produces the levelwise cones whose totalization is computed.
\item The product total complex of a double complex and its degreewise description. used at proof lines 30150-30151 and 30159-30165. Is the construction being compared, and the place where the degreewise product splitting occurs.
\item Acyclicity of a product totalization whose rows are acyclic. invoked at proof lines 30170-30172. Is the criterion that finishes the proof once the cone identification is made.
\end{itemize}\medskip
\ifdefined\hypertarget\hypertarget{sourcepointer-1}{}\fi
\noindent\textit{Source reference.} Original lemma 014D: focused-seeds/AGP/sources/stacks-snapshot-20261011/tags.txt, source lines 1468--1468. Cross-chapter navigation pointer into the Stacks Project; the standard fact it supplies is applied inside the counted proof span and the argument resting on it stays there.\par\medskip
\ifdefined\hypertarget\hypertarget{sourcepointer-2}{}\fi
\noindent\textit{Source reference.} Original lemma 30096: sources/stacks-more-algebra/pinned-original-source.tex, source lines 30096--30096. Same-chapter navigation pointer; the statement it names is standard for the section and the argument that consumes it stays inside the counted proof span.\par\medskip
\begin{lemma}
\label{lemma-prod-qis-gives-qis}
Let
$$
(A_0^\bullet \to A_1^\bullet \to A_2^\bullet \to \ldots)
\longrightarrow
(B_0^\bullet \to B_1^\bullet \to B_2^\bullet \to \ldots)
$$
be a map between two complexes of complexes of abelian groups.
Set $A^{p, q} = A_p^q$, $B^{p, q} = B_p^q$ to obtain double complexes.
Let $\text{Tot}_\pi(A^{\bullet, \bullet})$
and $\text{Tot}_\pi(B^{\bullet, \bullet})$ be the
product total complexes associated to the double complexes.
If each $A_p^\bullet \to B_p^\bullet$ is a
quasi-isomorphism, then
$\text{Tot}_\pi(A^{\bullet, \bullet}) \to \text{Tot}_\pi(B^{\bullet, \bullet})$
is a quasi-isomorphism.
\end{lemma}

\begin{proof}
Recall that $\text{Tot}_\pi(A^{\bullet, \bullet})$ in degree
$n$ is given by $\prod_{p + q = n} A^{p, q} = \prod_{p + 1 = n} A^q_p$.
Let $C_p^\bullet$ be the cone on the map $A_p^\bullet \to B_p^\bullet$,
see Derived Categories, Section \ref{derived-section-cones}.
By the functoriality of the cone construction we obtain a
complex of complexes
$$
C_0^\bullet \to C_1^\bullet \to C_2^\bullet \to \ldots
$$
Then we see $\text{Tot}_\pi(C^{\bullet, \bullet})$ in degree $n$
is given by
$$
\prod_{p + q = n} C^{p, q} = \prod_{p + q = n} C^q_p =
\prod_{p + q = n} (B^q_p \oplus A^{q + 1}_p) =
\prod_{p + q = n} B^q_p \oplus \prod_{p + q = n} A^{q + 1}_p
$$
We conclude that $\text{Tot}_\pi(C^{\bullet, \bullet})$
is the cone of the map
$\text{Tot}_\pi(A^{\bullet, \bullet}) \to \text{Tot}_\pi(B^{\bullet, \bullet})$
(We omit the verification that the differentials agree.)
Thus it suffices to show $\text{Tot}_\pi(A^{\bullet, \bullet})$ is
acyclic if each $A_p^\bullet$ is acyclic. This follows from
Lemma \ref{lemma-vanishing-coh-prod-totalization}.
\end{proof}

\end{document}
